% Front matter — the preface.  Short by design: the introduction owns
% the problem, the thesis, the reader contract, the data statement and
% the map; the preface says only what a reader should know before the
% table of contents.  No acknowledgements are invented here; the author
% adds them if and when he wants them.
\chapter*{Preface}
\markboth{Preface}{Preface}
\addcontentsline{toc}{chapter}{Preface}

This book is about publishing implied-volatility surfaces from listed
option quotes: models in which every fitted curve is a probability
law, the preparation of the quotes those models consume, and
inference for the mornings when the quotes run short.  It was written
alongside a working system --- the \emph{reference implementation}
that the text names throughout --- and that origin sets its
character.  Every construction in the book runs; every figure and
every empirical number regenerates from one frozen snapshot of market
data; and wherever mathematics ends and choice begins, the choice is
stated as one.

The book reads in order.  Each chapter assumes the ones before it,
read once, and the examples compound deliberately: the December smile
fitted in Part~I is still at work in the final chapter, where a
solver that never saw it reconstructs it from its neighbours.  The
introduction states the problem, the thesis that organizes every
chapter, what is assumed of the reader, and the map of what follows.

\vspace{2\baselineskip}
\begin{flushright}
Thibaut Delahaye\\
August 2026
\end{flushright}
